% GENERATED FILE. Edit canonical lessons, then run npm run generate:books.
% canonical-source-hash: fnv1a64:857c3488fba47731
% canonical-lessons: HI-C266-citkani, HI-C266-ari, HI-C266-kudal, HI-C266-phavra, HI-C266-gadda

\chapter{A latch, A saw, A pickaxe, A spade, A mattress}
\label{ch:hi-a-latch-a-saw-a-pickaxe-a-spade-a-mattress}
\hlchaptermodality{266}

\begin{chapteropening}
\textbf{By the end of this chapter:} \emph{I can say a latch, a saw, a pickaxe, a spade and a mattress.}

Complete the last lesson of chapter 266: i can say a latch, a saw, a pickaxe, a spade and a mattress.
\end{chapteropening}

\section[\texorpdfstring{\hi{चिटकनी} (ciṭkanī)}{चिटकनी (ciṭkanī)}]{\hi{चिटकनी} (ciṭkanī) — a latch}

\label{lesson:HI-C266-citkani}



\begin{quote}
Before the new one: say the Hindi for a cotton rug, then the Hindi for a carpet.
\end{quote}

\subsection*{You'll want to know: \hi{चिटकनी}}
\textbf{\hi{चिटकनी}} — \emph{ciṭkanī} — \textquotedblleft{}a latch\textquotedblright{}.

\begin{grammarlens}[title={What you've built}]
One more word for this chapter.
\end{grammarlens}

\subsection*{Your turn}
\begin{itemize}
\raggedright
  \item \emph{Say it:} \emph{ciṭkanī}
  \item \emph{Say it:} \emph{ciṭkanī}, once more
  \item \emph{Say it:} say \emph{ciṭkanī}
  \item \emph{Recall:} say the Hindi for a cotton rug, then the Hindi for a carpet, then say \emph{ciṭkanī} again
\end{itemize}

\begin{tcolorbox}[breakable,colback=teal!4,colframe=teal!35!black,title={Before you move on}]
What does \hi{चिटकनी} mean? (\textquotedblleft{}A latch\textquotedblright{}.) Say it once more.
\end{tcolorbox}
\section[\texorpdfstring{\hi{आरी} (ārī)}{आरी (ārī)}]{\hi{आरी} (ārī) — a saw}

\label{lesson:HI-C266-ari}



\begin{quote}
Before the new one: say the Hindi for a carpet, then the Hindi for a latch.
\end{quote}

\subsection*{You'll want to know: \hi{आरी}}
\textbf{\hi{आरी}} — \emph{ārī} — \textquotedblleft{}a saw\textquotedblright{}.

\begin{grammarlens}[title={What you've built}]
One more word for this chapter.
\end{grammarlens}

\subsection*{Your turn}
\begin{itemize}
\raggedright
  \item \emph{Say it:} \emph{ārī}
  \item \emph{Say it:} \emph{ārī}, once more
  \item \emph{Say it:} say \emph{ārī}
  \item \emph{Recall:} say the Hindi for a carpet, then the Hindi for a latch, then say \emph{ārī} again
\end{itemize}

\begin{tcolorbox}[breakable,colback=teal!4,colframe=teal!35!black,title={Before you move on}]
What does \hi{आरी} mean? (\textquotedblleft{}A saw\textquotedblright{}.) Say it once more.
\end{tcolorbox}
\section[\texorpdfstring{\hi{कुदाल} (kudāl)}{कुदाल (kudāl)}]{\hi{कुदाल} (kudāl) — a pickaxe, a hoe}

\label{lesson:HI-C266-kudal}



\begin{quote}
Before the new one: say the Hindi for a latch, then the Hindi for a saw.
\end{quote}

\subsection*{You'll want to know: \hi{कुदाल}}
\textbf{\hi{कुदाल}} — \emph{kudāl} — \textquotedblleft{}a pickaxe, a hoe\textquotedblright{}.

\begin{grammarlens}[title={What you've built}]
One more word for this chapter.
\end{grammarlens}

\subsection*{Your turn}
\begin{itemize}
\raggedright
  \item \emph{Say it:} \emph{kudāl}
  \item \emph{Say it:} \emph{kudāl}, once more
  \item \emph{Say it:} say \emph{kudāl}
  \item \emph{Recall:} say the Hindi for a latch, then the Hindi for a saw, then say \emph{kudāl} again
\end{itemize}

\begin{tcolorbox}[breakable,colback=teal!4,colframe=teal!35!black,title={Before you move on}]
What does \hi{कुदाल} mean? (\textquotedblleft{}A pickaxe, a hoe\textquotedblright{}.) Say it once more.
\end{tcolorbox}
\section[\texorpdfstring{\hi{फावड़ा} (phāvṛā)}{फावड़ा (phāvṛā)}]{\hi{फावड़ा} (phāvṛā) — a spade}

\label{lesson:HI-C266-phavra}



\begin{quote}
Before the new one: say the Hindi for a saw, then the Hindi for a pickaxe.
\end{quote}

\subsection*{You'll want to know: \hi{फावड़ा}}
\textbf{\hi{फावड़ा}} — \emph{phāvṛā} — \textquotedblleft{}a spade\textquotedblright{}.

\begin{grammarlens}[title={What you've built}]
One more word for this chapter.
\end{grammarlens}

\subsection*{Your turn}
\begin{itemize}
\raggedright
  \item \emph{Say it:} \emph{phāvṛā}
  \item \emph{Say it:} \emph{phāvṛā}, once more
  \item \emph{Say it:} say \emph{phāvṛā}
  \item \emph{Recall:} say the Hindi for a saw, then the Hindi for a pickaxe, then say \emph{phāvṛā} again
\end{itemize}

\begin{tcolorbox}[breakable,colback=teal!4,colframe=teal!35!black,title={Before you move on}]
What does \hi{फावड़ा} mean? (\textquotedblleft{}A spade\textquotedblright{}.) Say it once more.
\end{tcolorbox}
\section[\texorpdfstring{\hi{गद्दा} (gaddā)}{गद्दा (gaddā)}]{\hi{गद्दा} (gaddā) — a mattress}

\label{lesson:HI-C266-gadda}



\begin{quote}
Before the new one: say the Hindi for a pickaxe, then the Hindi for a spade.
\end{quote}

\subsection*{You'll want to know: \hi{गद्दा}}
\textbf{\hi{गद्दा}} — \emph{gaddā} — \textquotedblleft{}a mattress\textquotedblright{}.

\begin{grammarlens}[title={What you've built}]
One more word for this chapter.
\end{grammarlens}

\subsection*{Your turn}
\begin{itemize}
\raggedright
  \item \emph{Say it:} \emph{gaddā}
  \item \emph{Say it:} \emph{gaddā}, once more
  \item \emph{Say it:} say \emph{gaddā}
  \item \emph{Recall:} say the Hindi for a pickaxe, then the Hindi for a spade, then say \emph{gaddā} again
\end{itemize}

\begin{tcolorbox}[breakable,colback=teal!4,colframe=teal!35!black,title={Before you move on}]
What does \hi{गद्दा} mean? (\textquotedblleft{}A mattress\textquotedblright{}.) Say it once more.
\end{tcolorbox}
